% mayanumber-doc.tex — package documentation
% Compile twice with: xelatex mayanumber-doc.tex
\documentclass[11pt]{article}

\usepackage[a4paper, margin=2.5cm]{geometry}
\usepackage{mayanumber}      % the package under documentation
\usepackage{hyperref}
\hypersetup{colorlinks=true, linkcolor=mayaAgua, urlcolor=mayaAgua}

\title{\textbf{mayanumber}\\[4pt]
  \large Mayan numerals and vigesimal notation\\ with positional floor colors\\[6pt]
  version 1.0 --- 2026-09-09}
\author{Manuel López Mateos\\ \url{https://github.com/mlmateos}}
\date{}

\begin{document}
\maketitle

\begin{abstract}\noindent
The \texttt{mayanumber} package converts decimal integers of \emph{any}
magnitude into (a)~Mayan numerals (dots and bars, Unicode block
U+1D2E0--U+1D2F3) and (b)~vigesimal alphabetic notation (0--9, A--J),
coloring each positional \emph{floor} with a cycle of seven colors:
earth, water, vegetation, life, sky, sun, universe.
\end{abstract}

\section{Requirements and installation}

XeLaTeX or LuaLaTeX. The package loads \texttt{xcolor} and \texttt{fontspec}
by itself. It requires the free font \emph{BabelStone Mayan Numerals}
(\url{https://www.babelstone.co.uk/Fonts/}). To use another face, declare it
\emph{before} loading the package:

\begin{verbatim}
\newfontface{\mayanumerals}{Noto Sans Mayan Numerals}
\usepackage{mayanumber}
\end{verbatim}

Installation: drop \texttt{mayanumber.sty} next to your document, or into
\texttt{tex/latex/mayanumber/} of your local TEXMF tree.

\section{The command and its options}

\begin{verbatim}
\mayanumber{4581}                    % vertical + colors (default)
\mayanumber[horizontal]{4581}        % floors in a row, MSB left
\mayanumber[nocolor]{4581}           % plain glyphs
\mayanumber[nocolor, scale=3]{4581}  % options combine freely
\mayanumber[align=t]{4581}           % vertical alignment: c, t, b
\mayanumber*{4581}                   % alphabetic: B91 with subscript 20
\end{verbatim}

\begin{center}
\begin{tabular}{lllp{6cm}}
\textbf{Key} & \textbf{Values} & \textbf{Default} & \textbf{Meaning}\\\hline
\texttt{vertical} / \texttt{horizontal} & -- & vertical & stack floors, or lay them in a row\\
\texttt{colors} / \texttt{nocolor} & -- & colors & seven-color cycle on/off\\
\texttt{scale} & number & 1 & font scale\\
\texttt{align} & c, t, b & c & vertical alignment of the stack\\
\end{tabular}
\end{center}

Every call resets to the defaults before applying its own options, so no
state leaks between calls.

\section{The seven-floor color cycle}

Floors are colored from the bottom up (and right to left in horizontal mode):

\begin{center}
	\begin{tabular}{rlll}
		\textbf{Floor} & \textbf{Swatch} & \textbf{Name} & \textbf{Meaning}\\\hline
		6 & \colorbox{mayaUniverso}{\strut\hspace{1em}} & \#6A0DAD & the universe\\
		5 & \colorbox{mayaSol}{\strut\hspace{1em}} & \#DAA520 & light, the sun\\
		4 & \colorbox{mayaCielo}{\strut\hspace{1em}} & \#87CEEB & the sky\\
		3 & \colorbox{mayaVida}{\strut\hspace{1em}} & \#DC143C & life\\
		2 & \colorbox{mayaVegetacion}{\strut\hspace{1em}} & \#2E8B57 & vegetation\\
		1 & \colorbox{mayaAgua}{\strut\hspace{1em}} & \#005B96 & water\\
		0 & \colorbox{mayaTierra}{\strut\hspace{1em}} & \#8B4513 & the earth\\
	\end{tabular}\\[4pt]
	\emph{The table reads like a Mayan stack: the earth below, the universe above.}
\end{center}

\noindent Beyond floor~6 the cycle repeats: the universe becomes earth again.
The two blues differ in luminosity on purpose: water is deep and saturated,
sky is light and airy. The colors are declared with \verb|\providecolor|
(\texttt{mayaTierra}, \texttt{mayaAgua}, \dots), so users may redefine them
before loading the package.

\section{Examples}

\begin{center}
\mayanumber{4581}\qquad
\mayanumber[horizontal]{4581}\qquad
\mayanumber[nocolor]{4581}\qquad
\mayanumber*{4581}
\end{center}

\begin{center}
Seven floors, full cycle: \mayanumber[horizontal]{1234567890}\\[6pt]
Beyond 32 bits: \mayanumber[horizontal]{2187624240}
\end{center}

\section{How it works}

The decimal input is kept as a \emph{string} and divided by~20 with the
schoolbook long-division algorithm, digit by digit; the final remainder is
the floor digit and the quotient recurses until exhausted. All intermediate
values stay below~210, so TeX's 32-bit integer limit never applies and the
magnitude is unbounded. The color index of floor~$p$ is
$((p \bmod 7)+7) \bmod 7 + 1$, a residue guaranteed positive.

\section{Credits and license}

Original macro by \textbf{Enrico Gregorio (egreg on TeX StackExchange, @eg9 on GitHub)}, TeX StackExchange, in answer to a question
by \textbf{tatojo}: \url{https://tex.stackexchange.com/a/452806/18280}.
Extensions (seven-color cycle, vertical+color defaults, \texttt{nocolor},
unlimited magnitude) by Manuel López Mateos, 2026.
License: MIT; see the accompanying \texttt{LICENSE}.

\section{History}

\begin{tabular}{ll}
1.0 & 2026-09-09 --- first public release (GitHub and CTAN).\\
\end{tabular}

\end{document}
